% !TeX root = ../thesis.tex
% ===================== ГЛАВА 4 =====================
\chapter{Программная реализация и расчёты на суперкомпьютере}

\section{Прототип солвера}

Реализованный солвер представляет собой единую кодовую базу, которая отрабатывает численную схему и проводит тестовые
расчёты задач релаксации и переноса.
Текущая реализация~--- исследовательский прототип, предназначенный для отработки и верификации; вопрос промышленного
высокопроизводительного солвера выделен в отдельное направление (см.~\ref{sec:hpcpath}).

\section{Оптимизации, продиктованные размером задачи}

Размер фазовой сетки (порядка $N^3$ узлов по скоростям на каждую из $N_x$ пространственных ячеек) делает ключевыми две
оптимизации.
Во-первых, геометрия столкновений (узлы проекции, веса $r_\nu$, модули относительных скоростей) \emph{предвычисляется}
один раз, после чего вклад в интеграл столкновений набирается векторно по всем ячейкам, что даёт ускорение примерно на
два порядка по сравнению со скалярной реализацией.
Во-вторых, симметрия функции распределения позволяет хранить её «четвертью» (или «восьмой частью») массива, снижая
расход памяти в $4$--$8$ раз.

\section{Расчёты на суперкомпьютерном кластере}

Работа была выполнена с использованием оборудования центра коллективного пользования «Комплекс моделирования и обработки
данных исследовательских установок мега-класса» НИЦ «Курчатовский институт», http://ckp.nrcki.ru/.
Запуск, контроль состояния, просмотр журналов и выгрузка результатов автоматизированы (цикл
\emph{submit\,/\,status\,/\,log\,/\,fetch}).
Для независимых конфигураций использован механизм \emph{job-array}~--- одновременный счёт нескольких задач (например,
трёх значений параметра $u$ в задаче о квази-шоке).
Так, окончательный расчёт задачи о поршне ($N_x=400$, $p=50021$, $2000$ шагов) потребовал порядка $6 \cdot 10^{9}$
вычислений в узлах сетки Коробова.
Такой длительный прогон, необходимый для выхода ударной волны на стационар, на персональной рабочей станции практически
недостижим.

\section{Структурная параллелизуемость}

Алгоритм по своей природе хорошо распараллеливается: интеграл столкновений вычисляется независимо в каждой
пространственной ячейке, а узлы сетки Коробова независимы между собой.
Это позволяет естественно распределять работу по узлам кластера без существенных коммуникаций между ними.
В однородных задачах релаксации распараллеливание тривиально: независимые конфигурации (например, три значения $u$ в
задаче о квази-шоке) считаются одновременно механизмом \emph{job-array}.
В неоднородной задаче переноса естественная декомпозиция~--- по пространственным ячейкам: оператор столкновений в каждой
ячейке независим, а обмен между подобластями требуется только на шаге переноса и ограничен узкими слоями фиктивных ячеек
на границах.
Таким образом, объём межпроцессорных коммуникаций мал по сравнению с объёмом вычислений, что и обеспечивает хорошую
масштабируемость метода на суперкомпьютерных системах.

\section{Путь к промышленному солверу}
\label{sec:hpcpath}

Прямым продолжением работы является промышленная реализация солвера: компилируемый код (C/C++/Fortran) с
распараллеливанием по пространству средствами MPI и по узлам скоростной сетки на графических ускорителях (GPU)
